% Осень26!B31; нумерация внутри столбца без пустых строк.
% Источник: лекция 6, §4 (строка 198)
\definitionticket{28}{Образ меры при отображении}
\label{definition:28}

\begin{definition}
Пусть $(X,\mathcal A,\mu)$ --- пространство с мерой, $(Y,\mathcal B)$ --- множество с $\sigma$-алгеброй, а $\varphi:X\to Y$ --- измеримое отображение.
\textbf{Образом меры $\mu$ при отображении $\varphi$} называется мера $\nu$ на $\mathcal B$, заданная формулой
\[
  \nu(E):=\mu\bigl(\varphi^{-1}(E)\bigr),\qquad E\in\mathcal B.
\]
\end{definition}
